\section{Hybrid System Modeling}

%%%%%%%%%%%%%%%%%%%%%%%%%%%%%%%%%%%%%%%%%%%%%%%%%%%%%%%%%%%%%%%%%%%
\subsection{Hybrid System Representation}
\label{sec:hybrid_system_representation}

Cyber-physical systems combine continuous physical processes with digital control and software.
Further, impacts, switching logic, and sampled controllers can cause abrupt changes in modes, derivatives, or states.
These changes and transitions happen in a much faster way than the usual changes in the system and are therefore considered as instantaneous events.
A purely continuous model cannot describe such behavior in a consistent way.
Hybrid models therefore combine continuous dynamics with discrete change \cite{cellier_combined_1986,derler_modeling_2012}.

In this thesis, we use the reduced hybrid representation
\begin{equation}
    \left(q(t), x(t), z(t), u(t), y(t)\right),
\end{equation}
where $q(t)$ is the discrete state, $x(t)$ is the continuous state, $z(t)$ collects algebraic variables, $u(t)$ is the input, and $y(t)$ is the output.
The discrete state $q$ may represent an operating mode, a Boolean logic state, or an internal controller state.
This notation is a reduced representation inspired by formal hybrid modeling approaches such as Omola \cite{andersson_omola_1990,andersson_object-oriented_1994}.

Between events, the discrete state is constant and the continuous part follows the active continuous-time model:
\begin{equation}
    0 = F_{q(t)}\!\left(\dot{x}(t), x(t), z(t), u(t), t\right),
    \qquad
    y(t) = h_{q(t)}\!\left(x(t), z(t), u(t), t\right),
    \qquad t \neq t_e .
\end{equation}
At an event time $t_e$, the model may change its discrete state and may reset selected continuous states:
\begin{equation}
    x^{+} = R_j\!\left(x^{-}, z^{-}, q^{-}, u^{-}, t_e\right),
    \qquad
    q^{+} = \delta_j\!\left(x^{-}, z^{-}, q^{-}, u^{-}, t_e\right).
\end{equation}
The function $R_j(\cdot)$ computes the post-event values of the continuous states.
The function $\delta_j(\cdot)$ updates the discrete state.
Here, $x^{-}$ and $q^{-}$ denote the values immediately before the event, and $x^{+}$ and $q^{+}$ denote the values immediately after the event.
The system trajectory is therefore piecewise continuous \cite{van_der_schaft_introduction_2000,branicky_general_1995}.

Two event effects are important in this thesis.
A switching event changes the active mode or equations, but the continuous state remains continuous.
A reset event changes one or more continuous states instantaneously.
The first changes the derivative, while the second produces a jump in the state trajectory \cite{branicky_general_1995}.

A simple example is an integrator with reset:
\begin{equation}
    \dot{x}(t) = 1,
    \qquad
    y(t) = x(t).
\end{equation}
When the state reaches a threshold, an event is triggered and the state is reset:
\begin{equation}
    x \geq 1
    \quad \Rightarrow \quad
    x^{+} = 0.
\end{equation}
The state evolves continuously until the threshold is reached.
At the event time, the state jumps to zero.
After the reset, continuous evolution starts again from the new state.

%%%%%%%%%%%%%%%%%%%%%%%%%%%%%%%%%%%%%%%%%%%%%%%%%%%%%%%%%%%%%%%%%%%
\subsection{Events, Zero-Crossings, and State Re-Initialization}

Hybrid models change their behavior at event times.
Two event trigger types are common.
Time events occur at known or scheduled instants.
State events depend on the continuous trajectory and are therefore not known in advance \cite{cellier_combined_1986,branicky_general_1995}.

In this thesis, state events are represented by event indicators, also called zero-crossing functions,
\begin{equation}
    g_j\!\left(x(t), z(t), q(t), u(t), t\right).
\end{equation}
An event of type $j$ is triggered when the indicator crosses zero.
For example, the condition $x \geq x_{\max}$ can be written as
\begin{equation}
    g(x) = x - x_{\max}.
\end{equation}
The event occurs when $g(x)=0$ \cite{branicky_simulation_2006}.

In numerical simulation, the event time is usually not located exactly at the end of an integration step.
The simulator first detects that a zero-crossing occurred within a time interval.
It then localizes the event time inside this interval with a root-finding procedure.
This avoids integrating across a discontinuity without resolving its time of occurrence \cite{andersson_omola_1990,cellier_combined_1986}.

Once the event time $t_e$ has been found, the update maps introduced in Section~\ref{sec:hybrid_system_representation} are evaluated.
The resulting post-event state must satisfy the active continuous and algebraic equations before simulation can continue.
This step is called state re-initialization or restart \cite{andersson_omola_1990,urquia_moraleda_modeling_2018}.

%%%%%%%%%%%%%%%%%%%%%%%%%%%%%%%%%%%%%%%%%%%%%%%%%%%%%%%%%%%%%%%%%%%
\subsection{Time Semantics in Hybrid Simulation}

Hybrid simulation alternates between phases of continuous evolution and discrete event handling.
During continuous phases, the state evolves according to the active differential-algebraic equations.
At event times, discrete updates may change the discrete state, reset continuous states, or trigger further events.

A purely real-valued notion of time is often not sufficient for this process.
Several discrete updates may have to be executed at the same physical time.
This occurs, for example, when one event triggers another event immediately, or when several event conditions become active at the same instant \cite{urquia_moraleda_modeling_2018}.

For this reason, hybrid simulation often uses \emph{superdense time}.
A time point is then represented by the pair
\begin{equation}
    (t,n),
\end{equation}
where $t \in \mathbb{R}_{\geq 0}$ denotes the physical time and $n \in \mathbb{N}_0$ denotes a logical micro-step index.
The physical time identifies when the event occurs.
The index $n$ orders the sequence of discrete updates at that same physical time.
An increase in $n$ does not represent additional physical time.
It only represents progress in the event iteration at fixed time $t$ \cite{broman_determinate_2013,cremona_hybrid_2019}.

Superdense time makes it possible to distinguish between the state before an event and the states after one or more discrete updates at the same physical time.
Continuous integration resumes only after all event updates at the current physical time have been processed.
This distinction is important for the present thesis, since the hybrid co-simulation algorithm introduced later first localizes an event in physical time and then performs ordered event handling at that same time instant.